\begin{table}[ht]
\caption{\textbf{\pony example}}
\centering
% \resizebox{\textwidth}{!}{
\begin{tabular}{p{13cm}}
\toprule
% \hspace{5.5cm} \econ \\
% \midrule
\textit{\textbf{Query}} \\
\midrule
     Given the lengths of a triangle's sides, write a pony program to classify it as equilateral, isosceles or scalene.
 \\
\midrule
\textit{\textbf{Chain-of-thought reasoning to find documents}} \\
\midrule
To determine the category of the triangle, we need to compare the side lengths. A common syntax for conditions and comparison used in many other programming languages is the `if' or `switch' related to control structure. \\
We need to find the syntax that allows to specify such conditions. \\
\midrule
\textit{\textbf{Example positive document}} \\
\midrule
\# Control Structures \\
\\
To do real work in a program you have to be able to make decisions, iterate \\ through collections of items and perform actions repeatedly. For this, you need \\ control structures. Pony has control structures that will be familiar to programmers who have used most languages, \\
such as `if', `while' and `for', but in Pony, they work slightly differently. \\
\\
\#\# Conditionals \\
\\
The simplest control structure is the good old `if`. It allows you to perform some action only when a condition is true. \\
In Pony it looks like this: \\
\\
if condition then  \\
  control\_body \\
end \\
\\
Here is a simple example: \\
\\
if a > b then \\
  env.out.print("a is bigger") \\
end \\
\\
Often the condition may be composed of many sub conditions connected by `and' and `or'. \\
... \\
\midrule
\textit{\textbf{Example negative document}} \\
\midrule
\# Classes \\
\\
Just like other object-oriented languages, Pony has \_\_classes\_\_. \\
A class is declared with the keyword `class`, and it has to have a name \\
that starts with a capital letter, like this: \\
\\
class Wombat \\
\\
Do all types start with a capital letter?\\
Yes! And nothing else starts with a capital letter. \\
So when you see a name in Pony code, you will instantly know whether it's a type or not. \\
\\
\#\# What goes in a class? \\
\\
A class is composed of: \\
\\
1. Fields. \\
2. Constructors. \\
3. Functions. \\
% \\
% \#\#\# Fields \\
% \\
% These are just like fields in C structures or fields in classes in C++, C\#, \\
% Java, Python, Ruby, or basically any language, really. There are three kinds of fields: \\
% `var', `let', and `embed' fields. A `var' field can be assigned to over and \\
% over again, but a `let` field is assigned to in the constructor and never \\
% again. Embed fields will be covered in more detail in the documentation on \\ variables. \\
% \\
% class Wombat \\
%   let name: String \\
%   var \_hunger\_level: U64 \\
% \\
% Here, a `Wombat` has a `name`, which is a `String`, and a `\_hunger\_level`, \\ which is a `U64` (an unsigned 64-bit integer).

% What does the leading underscore mean? \\
% It means something is \_\_private\_\_. \\
% A \_\_private\_\_ field can only be accessed by code in the same type. \\
% A \_\_private\_\_ constructor, function, or behaviour can only be accessed \\
% by code in the same package. We'll talk more about packages later. \\
... \\
\bottomrule
\end{tabular}
% }
\label{tab:pony_example}
\end{table}